% ==============================================================================
% CAPÍTULO 4: RESULTADOS
% O título muda conforme a etapa (\titresultados, definido em main.tex):
% "Resultados preliminares" na qualificação e "Resultados e discussão" na Defesa.
% ==============================================================================
\chapter{\titresultados}
\label{ch:analise_discussao}
\begin{resumocapitulo}
  Aqui deve ser incluído um resumo do capítulo, descrevendo, de forma reduzida, os achados até o momento, de acordo com os objetivos propostos com o limite de 200 palavras.
\end{resumocapitulo}

\preencher{A análise e a discussão dos resultados preliminares têm como finalidade examinar e contextualizar criticamente os primeiros achados empíricos obtidos na investigação, confrontando-os com o quadro teórico previamente estabelecido. O seu papel estrutural é interpretar o comportamento das variáveis, padrões emergentes ou evidências parciais à luz da literatura e dos objetivos da pesquisa, avaliando a consistência dos modelos analíticos aplicados. Essa seção permite diagnosticar a adequação dos protocolos metodológicos em curso, subsidiar eventuais ajustes analíticos e antecipar as tendências que orientarão a consolidação das conclusões finais da dissertação. Apagar este texto de instrução após a redação da seção.}

\section{Visão geral}
\label{sec:visao_geral_resultados}
Descrever uma síntese do que será tratado no capítulo, incluindo os objetivos específicos e os resultados preliminares obtidos. A visão geral deve fornecer ao leitor um panorama dos achados, destacando como cada resultado contribui para responder à pergunta de pesquisa e alcançar os objetivos propostos. Podem ser utilizadas outras subseções para organizar os resultados por objetivo específico, por construto teórico ou por tipo de análise realizada.

Exemplo de inserção de tabela com os resultados numéricos. A Tabela~\ref{tab:exemplo} apresenta resultados numéricos aleatórios. Substituir com seus dados.
\begin{table}[H]
  \centering
  \caption{Exemplo de tabela com apresentação de dados numéricos}
  \label{tab:exemplo}
  \begin{tabular*}
    {\linewidth}
    {@{\extracolsep{\fill}} l c}
    \hline \textbf{Métrica} & \textbf{Valor} \\ \hline Acurácia (\textit{accuracy}) & 0,51 \\ Precisão (\textit{precision}) & 0,50 \\ Revocação (\textit{recall}) & 0,96 \\ \textit{F1-score} & 0,66 \\ \hline
  \end{tabular*}
  \fonte{Elaborada pelo autor (\ano).}
\end{table}
